\chapter{Практика}


\section{Настройка интерфейсов роутера}

\begin{lstlisting}[language=C, caption={Пример настройки интерфейсов роутера}]
#R1
configure

interface gigabitethernet 1/0/1
ip address 192.168.100.1/30
ip firewall disable
exit

interface gigabitethernet 1/0/2
ip address 10.10.0.1/24
ip firewall disable
exit

1interface gigabitethernet 1/0/3
ip address 10.10.1.1/24
ip firewall disable
exit

interface gigabitethernet 1/0/4
ip address 10.10.100.1/30
ip firewall disable
exit

do commit
do confirm

#R2
configure

interface gigabitethernet 1/0/1
ip address 192.168.100.2/30
ip firewall disable
exit

interface gigabitethernet 1/0/2
ip address 10.10.10.1/24
ip firewall disable
exit

1interface gigabitethernet 1/0/3
ip address 10.10.11.1/24
ip firewall disable
exit

interface gigabitethernet 1/0/4
ip address 172.16.200.2/30
ip firewall disable
exit

do commit
do confirm

#R3
configure

interface gigabitethernet 1/0/1
ip address 172.16.200.1/30
ip firewall disable
exit

interface gigabitethernet 1/0/2
ip address 10.10.100.2/30
ip firewall disable
exit

do commit
do confirm
\end{lstlisting}

\section{Настройка ip на ПК}

\begin{lstlisting}[language=C, caption={Пример настройки ip на ПК}]
#VPC1
ip 172.16.10.10/24 172.16.10.1
#VPC2
ip 172.16.20.20/24 172.16.20.1
#VPC3
ip 172.16.30.30/24 172.16.30.1
#VPC4
ip 10.10.40.10/24 10.10.40.1
#VPC5
ip 10.10.50.20/24 10.10.50.1
#VPC6
ip 10.10.60.30/24 10.10.60.1
\end{lstlisting}

\section{Настройка главных статических маршрутов}

\begin{lstlisting}[language=C, caption={Пример настройки статических маршрутов}]
# R1-R2
ip route 10.10.10.0/23 192.168.100.2

do commit
do confirm 

#R2-R1
ip route 10.10.0.0/23 192.168.100.1

do commit
do confirm 
\end{lstlisting}

Используем маску /23, т.к она охватывает все подсети на каждом из маршрутизаторов.

\section{Настройка плавающих статических маршрутов}

\begin{lstlisting}[language=C, caption={Пример настройки плавающих маршрутов}]
# R1-R3
ip route 10.10.10.0/23 10.10.100.2 10

do commit
do confirm 

#R2-R3
ip route 10.10.0.0/23 172.16.200.1 10

do commit
do confirm 
\end{lstlisting}

10 в конце означает cost. Мы повышаем "стоимость" пути, чтобы он хранился как запасной.

\section{Настройка путей для каждой подсети}

\begin{lstlisting}[language=C, caption={Пример настройки плавающих маршрутов}]
# R1-R2-R3
Router-down(config)# ip route 10.10.10.0/24 172.16.200.2
Router-down(config)# ip route 10.10.11.0/24 172.16.200.2

# R3-R2-R1
Router-down(config)# ip route 10.10.0.0/24 10.10.100.1
Router-down(config)# ip route 10.10.1.0/24 10.10.100.1

\end{lstlisting}

\section{Проверка}

Сделаем trace ПК1-ПК3

\begin{figure}[H] 
    \centering
    \includegraphics[width=0.75\textwidth]{t1.png}
    \caption{Прямой маршрут}
\end{figure}

Удалим соединение R1-R2 и повторим trace

\begin{figure}[H] 
    \centering
    \includegraphics[width=0.75\textwidth]{t2.png}
    \caption{Плавающий маршрут}
\end{figure}

Маршруты пакетов отличаются: в первом случае трафик шел через R1-R2, а во втором случае через R1-R2-R3, т.е маршрут стал актулальным после того, как пропал главный маршрут.

\endinput